\chapter*{Nota sobre este borrador}
\addcontentsline{toc}{chapter}{Nota sobre este borrador}

Este documento es un \textbf{borrador parcial} de la memoria del Trabajo Fin de Grado, preparado para una primera revisión por parte del tutor antes de la entrega definitiva (22~de~mayo de 2026). Su objetivo es validar la estructura, el enfoque y el alcance del proyecto.

\section*{Capítulos incluidos en este borrador}

\begin{itemize}
  \item \textbf{Resumen} y \textbf{abstract} en español e inglés.
  \item \textbf{Capítulo~1 — Introducción}: motivación, objetivos formales y declaración explícita del uso de IA.
  \item \textbf{Capítulo~2 — Estado del arte}: análisis de cinco aplicaciones de referencia, comparativa funcional y comparativa técnica de las tecnologías evaluadas.
  \item \textbf{Capítulo~3 — Análisis y diseño}: requisitos funcionales (46 RF) y no funcionales (12 RNF), actores, casos de uso, arquitectura, modelo de datos (13~entidades), flujo de estados y diseño de la API REST.
  \item \textbf{Capítulo~6 — Conclusiones}: logros, limitaciones honestas y trabajo futuro.
\end{itemize}

\section*{Capítulos en redacción (no incluidos aquí)}

\begin{itemize}
  \item \textbf{Capítulo~4 — Implementación}: desarrollo por sprints, decisiones técnicas, extractos de código. Listo en la versión completa, omitido aquí para no saturar la primera revisión.
  \item \textbf{Capítulo~5 — Evaluación}: resultados de pruebas, validación de requisitos. Pendiente de incorporar las pruebas de usabilidad con tres usuarios reales.
  \item \textbf{Anexos}: guía de instalación, manual de usuario y modelo de negocio.
\end{itemize}

\section*{Pendiente de incorporar antes de la entrega final}

\begin{itemize}
  \item Capturas de pantalla reales de la aplicación (8 capturas previstas).
  \item Fichas detalladas de las cinco apps del estado del arte tras pruebas personales.
  \item Resultados cuantitativos de las pruebas de usabilidad.
  \item Texto final de los agradecimientos.
\end{itemize}

\section*{Cuestiones para la próxima tutoría}

\begin{itemize}
  \item Confirmación del estilo de citas (la plantilla usa APA~7).
  \item Confirmación de la modalidad de defensa.
  \item Decisión sobre publicación en RUA.
  \item Cualquier observación sobre alcance, estructura o correcciones a aplicar antes de la entrega.
\end{itemize}

\noindent\rule{\textwidth}{0.4pt}
\begin{flushright}
\small Erardo Aldana Pessoa\\
Convocatoria C3 — Curso 2025--2026
\end{flushright}

\clearpage
